% !PRINTED 721
\weberpage{721}

Die Zahl $1 + \mu$ hat die Indices $\alpha = 0$, $\beta = \nu'$, und
folglich ist nach (7)
\[
\chi(1+\mu) = \Theta^{\nu'} = -1,
\]
wenn $\chi$ ein primitiver Charakter ist.

Ferner ist f\"ur ein ungerades $n$:
\[
n(1+\mu) \equiv n + \mu \pmod{m},
\]
und folglich
\[
\chi(n+\mu) = -\chi(n).
\]
Wenn wir nun die Function $f(r^s)$ betrachten:
\weq{10}{f(r^s) = \overset{n}{\Sigma} \chi(n)\, r^{sn},}
worin $n$ die ungeraden Zahlen eines vollen Restsystems f\"ur den Modul
$m$ durchl\"auft, so erhalten wir, da wir in (10) unter dem
Summenzeichen $n$ durch $n+\mu$ ersetzen d\"urfen:
\[
f(r^s) = -r^{\mu s} \sum \chi(n)\, r^{sn},
\]
und da $r^\mu = -1$ ist:
\[
f(r^s) = -(-1)^s f(r^s).
\]
Es ist also
\weq{11}{f(r^s) = 0, \quad \text{\gesperrt{wenn $s$ gerade ist,}}}
\gesperrt{und demnach k\"onnen wir in den Formeln (9) den Summationsbuchstaben $s$
auf die ungeraden Zahlen beschr\"anken, die kleiner als $m$ sind.}

Es ist aber nach (10):
\[
f(r^s) = \chi(s)^{-1} \overset{n}{\Sigma} \chi(ns)\, r^{sn},
\]
und wenn man $sn$ durch $n$ ersetzt, was bei ungeradem $s$ gestattet
ist:
\weq{12}{f(r^s) = \chi(s)^{-1} f(r), \quad \text{\gesperrt{wenn $s$ ungerade ist.}}}
Die Function $f(r)$ hat, wenn $\alpha, \beta$ die Indices von $n$ sind,
den Ausdruck
\weq{13}{f(r) = \overset{n}{\Sigma} (\pm 1)^\alpha \Theta^\beta r^n,}
und ist also mit der im \S.\,17 dieses Bandes betrachteten
Kreistheilungsresolvente
\[
(\pm 1, \Theta, r)
\]
gleichbedeutend, f\"ur die im \S.\,17, (17) die Relation aufgestellt
war:
\weq{14}{(\pm 1, \Theta, r)(\pm 1, \Theta^{-1}, r) = \pm m.}
Hier gelten durchweg die oberen Zeichen f\"ur die Summen $X_1$, die
unteren f\"ur die Summen $X_2$ [\S.\,195, (7)].

Wir theilen jetzt die Reihe der in (9) vorkommenden Zahlen $s$ in vier
Theile. Es soll $t$ ein Zeichen sein, welches die
%CONT>>
